% Generated draft: requires mathematical and source review.
\documentclass[11pt]{article}
\usepackage[margin=25mm]{geometry}
\usepackage{fontspec}
\setmainfont{Latin Modern Roman}
\usepackage{amsmath,amssymb,mathtools,unicode-math}
\setmathfont{Latin Modern Math}
\usepackage{enumitem,needspace,xurl,hyperref}
\setlength{\parindent}{0pt}
\setlength{\parskip}{5pt}
\setlength{\emergencystretch}{3em}
\newcommand{\N}{\mathbb{N}}
\newcommand{\Z}{\mathbb{Z}}
\newcommand{\Q}{\mathbb{Q}}
\newcommand{\R}{\mathbb{R}}
\newcommand{\C}{\mathbb{C}}
\newcommand{\F}{\mathbb{F}}
\newcommand{\A}{\mathbb{A}}
\newcommand{\PP}{\mathbb{P}}
\newcommand{\E}{\mathbb{E}}
\newcommand{\eps}{\varepsilon}
\newcommand{\dd}{\,\mathrm{d}}
\begin{document}
\section*{UnsolvedMath ID: 701528}
\textbf{Source problem code:} catalogue entry
\subsection{Definitions and mathematical statement}
For a finite group \(K\), let \(\operatorname{Irr}(K)\) denote the set of irreducible complex characters of \(K\), and define
\[
\operatorname{cd}(K)=\{\chi(1):\chi\in\operatorname{Irr}(K)\}.
\]
Thus \(\operatorname{cd}(K)\) is a set of positive integers, with multiplicities discarded.

\textbf{Huppert's Conjecture (supplied formulation).} For every finite nonabelian simple group \(H\) and every finite group \(G\),
\[
\operatorname{cd}(G)=\operatorname{cd}(H)
\quad\Longrightarrow\quad
\text{there exists a finite abelian group }A\text{ such that }G\cong H\times A.
\]
The group \(A\) may be trivial. The finiteness of \(H\) is implicit in the source's definition of character degrees for finite groups; it is made explicit here. The source calls \(H\) a nonabelian simple group without separately repeating this finiteness hypothesis.
\subsection{Short English statement}
Huppert's Conjecture asks whether every finite group with exactly the same set of irreducible complex character degrees as a finite nonabelian simple group \(H\) is a direct product of \(H\) and a finite abelian group.
\subsection{Sources}
\begin{itemize}
\item Supplied local UnsolvedMath registry record: permanent ID \(701528\), exact source problem code \texttt{catalogue entry}, title \emph{Kourovka Notebook Problem 18.99}; upstream revision \texttt{d371ee917ac9647fbb06cd6fc529deec6f3dda82}. The record was updated on \texttt{2026-10-16}.
\item The supplied attribution is \emph{The Kourovka Notebook (Unsolved Problems in Group Theory)}, 18th issue (2014), problem 18.99, proposed by H. P. Tong-Viet. The statement is named Huppert's Conjecture in the record.
\item The supplied background also identifies the 21st issue edition (2026), arXiv identifier \texttt{1401.0300v48}. The collection description names E. I. Khukhro and V. D. Mazurov as editors of Issues 1--21 (1965--2026). These bibliographic details are supplied attribution and have not been independently verified.
\item The registry labels this problem \texttt{open}. This is a supplied status label, not an independently verified assessment of the literature.
\end{itemize}
\end{document}
